\chapter{Testicular Pain (Orchalgia)}

\section*{Definition}
Testicular pain is discomfort in one or both testicles. It can be acute or chronic and may arise from the testicle itself, the epididymis, or referred from other structures. Testicular torsion is a urologic emergency that requires prompt surgical intervention.

\section*{Pathophysiology}
Testicular pain can be caused by inflammation (epididymitis, orchitis), torsion (twisting of the spermatic cord causing reduced blood flow), trauma, tumors, or varicocele (dilated veins). Torsion cuts off blood supply, causing reduced blood flow and infarction within 4-6 hours. Referred pain can occur from kidney stones, hernia, or prostatitis via shared nerve pathways.

\section*{Causes}
\begin{itemize}
\item Epididymitis (inflammation of the epididymis, often infectious)
\item Orchitis (inflammation of the testicle, often viral: mumps)
\item Testicular torsion (surgical emergency)
\item Torsion of the appendix testis
\item Varicocele (dilated pampiniform plexus)
\item Spermatocele (cystic dilation of the epididymis)
\item Trauma or blunt injury
\item Testicular tumor (usually painless but can cause dull ache)
\item Inguinal hernia
\item Referred pain from kidney stones or prostatitis
\end{itemize}

\section*{When to See a Doctor}
\begin{itemize}
\item Sudden, severe testicular pain (emergency - rule out torsion)
\item Dull ache or heavy sensation
\item Swelling or enlargement of the testicle
\item Redness or warmth of the scrotal skin
\item Nausea and vomiting (with torsion)
\item Fever (with epididymitis or orchitis)
\item Pain that improves with elevation (epididymitis)
\item Urethral discharge (with epididymitis)
\end{itemize}

\section*{Related Symptoms}
\begin{itemize}
\item Testicular swelling
\item Scrotal pain
\item Groin pain
\item Lower abdominal pain
\item Urinary symptoms (frequency, urgency, painful urination)
\end{itemize}

\section*{Self-Care / Home Management}
\begin{itemize}
\item Wear supportive underwear (briefs instead of boxers)
\item Apply ice packs to the scrotum
\item Elevate the scrotum
\item Take NSAIDs for pain
\item Avoid heavy lifting and strenuous activity
\item Do NOT apply heat (can worsen inflammation)
\item Seek immediate care for sudden, severe pain
\end{itemize}

\section*{How Your Doctor Will Evaluate You}
\begin{itemize}
\item Ultrasound with Doppler (urgent for suspected torsion)
\item Urinalysis and urine culture
\item Complete blood count
\item STI testing (chlamydia, gonorrhea) for epididymitis
\item Testicular tumor markers (AFP, hCG) if mass found
\end{itemize}

\section*{Treatment Options}
\begin{itemize}
\item Testicular torsion: immediate surgical exploration and orchiopexy
\item Epididymitis: antibiotics (doxycycline, ceftriaxone) and rest
\item Orchitis: supportive care, NSAIDs, scrotal elevation
\item Varicocele: observation or surgical ligation if symptomatic
\item Trauma: rest, ice, elevation, NSAIDs; surgery for rupture
\item Testicular tumor: radical orchiectomy and oncology referral
\end{itemize}
